\section*{Capítulo \ref{ch:g6:latitude-longitude} --- Latitude, longitude e fusos horários}

\begin{solution}{exo:g6:latitude-longitude:1}
``A \omterm{def:g6:latitude-longitude:latitude}{latitude} é medida a partir do \omterm{def:g3:continents-and-oceans:equator}{equador} e vai de $0^\circ$ a $90^\circ$
(norte ou sul); a \omterm{def:g6:latitude-longitude:longitude}{longitude} é medida a partir do \omterm{def:g6:latitude-longitude:longitude}{meridiano de origem} e vai
de $0^\circ$ a $180^\circ$ (leste ou oeste).''
\end{solution}

\begin{solution}{exo:g6:latitude-longitude:2}
O \omterm{def:g3:continents-and-oceans:poles}{Polo Norte}: $90^\circ$\,N. O \omterm{def:g3:continents-and-oceans:poles}{Polo Sul}: $90^\circ$\,S. Uma cidade no
\omterm{def:g3:continents-and-oceans:equator}{equador}: $0^\circ$.
\end{solution}

\begin{solution}{exo:g6:latitude-longitude:3}
Sydney: \omterm{def:g3:continents-and-oceans:equator}{hemisfério sul}, leste. Nova York: \omterm{def:g3:continents-and-oceans:equator}{hemisfério norte}, oeste. Cairo:
\omterm{def:g3:continents-and-oceans:equator}{hemisfério norte}, leste. Rio de Janeiro: \omterm{def:g3:continents-and-oceans:equator}{hemisfério sul}, oeste.
\end{solution}

\begin{solution}{exo:g6:latitude-longitude:4}
(a) Verdadeiro. (b) Falso: os \omterm{def:g6:latitude-longitude:latitude}{paralelos} ficam mais curtos na direção dos
\omterm{def:g3:continents-and-oceans:poles}{polos}; o mais longo é o \omterm{def:g3:continents-and-oceans:equator}{equador}. (c) Verdadeiro. (d) Falso: ele passa por
Greenwich, em Londres.
\end{solution}

\begin{solution}{exo:g6:latitude-longitude:5}
O Cairo. Nova York. Tóquio fica a cerca de $36^\circ$\,N, $140^\circ$\,L
(um pouco acima do \omterm{def:g6:latitude-longitude:latitude}{paralelo} $30^\circ$\,N, um pouco além do \omterm{def:g6:latitude-longitude:longitude}{meridiano}
$120^\circ$\,L, antes de $150^\circ$\,L).
\end{solution}

\begin{solution}{exo:g6:latitude-longitude:6}
$45 \div 15 = 3$: os seus meios-dias têm 3~horas de diferença. É meio-dia
primeiro na cidade~B, que fica mais a leste.
\end{solution}

\begin{solution}{exo:g6:latitude-longitude:7}
$73 \div 15 \approx 4.9$, arredondado para 5: o fuso é $+5$. Quando é
meio-dia em Greenwich, são $12 + 5 = 17$, ou seja, 17~h, em Mumbai.
\end{solution}

\begin{solution}{exo:g6:latitude-longitude:8}
$30 \times 111 = 3\,330$~km. O \omterm{def:g3:continents-and-oceans:poles}{Polo Norte} fica a $90^\circ$:
$90 \times 111 = 9\,990$~km, cerca de 10\,000~km.
\end{solution}

\begin{solution}{exo:g6:latitude-longitude:9}
$9 - (-5) = 14$~horas de diferença. $8 + 14 = 22$: são 22~h de
segunda-feira em Tóquio.
\end{solution}

\begin{solution}{exo:g6:latitude-longitude:10}
Ela tem razão. A Terra gira de oeste para leste, então as cidades do leste
são levadas na direção do Sol primeiro: elas o veem nascer, chegar ao seu
ponto mais alto e se pôr antes das cidades mais a oeste.
\end{solution}

\begin{solution}{exo:g6:latitude-longitude:11}
(a) $42 \times 3 = 126$~segundos, um pouco mais de 2~minutos. (b)
$126 \div 240 \approx 0.5$: cerca de meio grau de \omterm{def:g6:latitude-longitude:longitude}{longitude}. (c)
$0.5 \times 111 \approx 55$~km (com $0.525$: cerca de 58~km). Um erro de
cerca de cinquenta quilômetros: o bastante para não encontrar uma ilha, ou
para encalhar numa costa que se pensava estar longe.
\end{solution}

\begin{solution}{exo:g6:latitude-longitude:12}
(a) $18 - 9 = 9$: 9~h de segunda-feira em Greenwich. (b) $9 + 10 = 19$:
19~h de segunda-feira em Greenwich. (c) $19 - 8 = 11$: 11~h de
segunda-feira em Los Angeles. O avião saiu às 18~h de segunda-feira (hora
de Tóquio) e pousa às 11~h da mesma segunda-feira (hora de Los Angeles):
os relógios na chegada mostram uma hora mais cedo que os relógios na
partida, porque o avião atravessou a \omterm{def:g6:latitude-longitude:date-line}{linha de data} para o leste e voltou
um dia.
\end{solution}

\begin{solution}{pb:g6:latitude-longitude:1}
\textbf{1.} Ana (Rio de Janeiro, $23^\circ$\,S) e Mia (Auckland,
$37^\circ$\,S).

\textbf{2.} Ana e Sam ($43^\circ$\,O e $118^\circ$\,O). Lena está no
\omterm{def:g6:latitude-longitude:longitude}{meridiano de origem}.

\textbf{3.} Lena: Londres fica a $52^\circ$\,N, mais longe do \omterm{def:g3:continents-and-oceans:equator}{equador} que
todos os outros.

\textbf{4.} Mia: Auckland fica a $175^\circ$\,L, a só $180 - 175 = 5$
graus do \omterm{def:g6:latitude-longitude:longitude}{meridiano} de $180^\circ$.

\textbf{5.} $360^\circ \div 24 = 15^\circ$.

\textbf{6.} Londres: $0$. Rio de Janeiro: $43 \div 15 \approx 2.9$, oeste,
então $-3$. Tóquio: $140 \div 15 \approx 9.3$, leste, então $+9$. Auckland:
$175 \div 15 \approx 11.7$, leste, então $+12$. Los Angeles:
$118 \div 15 \approx 7.9$, oeste, então $-8$.

\textbf{7.} Rio de Janeiro: $20 - 3 = 17$, 17~h de sexta-feira. Tóquio:
$20 + 9 = 29$; $29 - 24 = 5$: 5~h de sábado.

\textbf{8.} Tóquio fica muito a leste: os seus relógios estão 9~horas
adiantados em relação aos de Londres. Às 20~h em Londres, Tóquio já passou
da meia-noite, então o seu sábado já começou.

\textbf{9.} $20 - 8 = 12$: meio-dia de sexta-feira.

\textbf{10.} $12 - (-8) = 20$~horas adiantados.

\textbf{11.} Atravessando a \omterm{def:g6:latitude-longitude:date-line}{linha de data} para o leste, o navio volta para
o dia anterior.

\textbf{12.} $20 + 12 = 32$; $32 - 24 = 8$: quando a chamada começa, são
\textbf{8~h de sábado} em Auckland. Mia entra na chamada na hora do café da
manhã, enquanto Sam almoça na sexta-feira.
\end{solution}
